%! Periodic Phenomena — Unit 6, Lesson 6.1 — Group Activity (Answer Key)
\documentclass[10pt]{article}
\usepackage{precalculus-article}
\usepackage{precalculus-key}   % pulls in -boxes; do NOT also load -boxes
\newcommand{\TallMath}[1]{$\displaystyle #1\rule[-1.4em]{0pt}{3.2em}$}

\begin{document}

\pageheader{Unit 6, Lesson 6.1}{Group Activity --- Answer Key}
\namepartnerperiod

\vspace{0.10in}

\begin{scenariobox}[Context: Periodic Patterns in the World]{plumlight}
Periodic relationships appear throughout nature, engineering, and everyday
life. In this activity your group will investigate several real-world
contexts, determine whether each is periodic, estimate periods, and describe
key characteristics. Choose \textbf{one tier} below based on your readiness.
\end{scenariobox}

\vspace{0.08in}

%----------------------------------------------------------------------
% Tier R Key
%----------------------------------------------------------------------
\begin{tcolorbox}[colback=white, colframe=black!40,
    title=\textbf{Tier R --- Reinforce (Key)},
    fonttitle=\bfseries, arc=2mm, left=3mm, right=3mm, top=2mm, bottom=2mm]

\textbf{Context:} Portland, OR average daily high temperature by month.

\begin{center}
\renewcommand{\arraystretch}{1.35}
\small
\begin{tabular}{|c|c|c|c|c|c|c|c|c|c|c|c|c|}
\hline
\textbf{Month} & 1 & 2 & 3 & 4 & 5 & 6 & 7 & 8 & 9 & 10 & 11 & 12 \\
\hline
\textbf{Year 1} & 47 & 52 & 58 & 63 & 70 & 76 & 82 & 82 & 74 & 63 & 52 & 47 \\
\hline
\textbf{Month} & 13 & 14 & 15 & 16 & 17 & 18 & 19 & 20 & 21 & 22 & 23 & 24 \\
\hline
\textbf{Year 2} & 47 & 52 & 58 & 63 & 70 & 76 & 82 & 82 & 74 & 63 & 52 & 47 \\
\hline
\end{tabular}
\end{center}

\vspace{6pt}
\textbf{1.}\quad Is this relationship periodic? Circle one: \quad \textbf{YES} \quad / \quad \textbf{NO}

Explain: \ans{Yes. The output values in Year 2 (months 13--24) are identical to those in Year 1 (months 1--12), repeating the same sequence over the same input length of 12 months.}

\vspace{4pt}
\textbf{2.}\quad Estimate the period. The output pattern restarts at month \ans{13},
so the period is $k =$ \ans{12} months.

\vspace{4pt}
\textbf{3.}\quad Describe one cycle:
\begin{itemize}[leftmargin=1.4em, itemsep=4pt]
  \item Increasing from month \ans{1} to month \ans{7 (or 8)}.
  \item Decreasing from month \ans{7 (or 8)} to month \ans{12}.
  \item Maximum: \ans{$82^\circ$F} at month \ans{7 and 8}.
  \item Minimum: \ans{$47^\circ$F} at month \ans{1 and 12}.
\end{itemize}

\vspace{4pt}
\textbf{4.}\quad \ans{Yes. Because the relationship is periodic with period 12, the output values in every subsequent period are identical to those in the first cycle. The same max and min will appear in months 25--36, 37--48, and so on.}

\end{tcolorbox}

\begin{teachernote}
\textbf{Common errors --- Tier R:}
\begin{itemize}[leftmargin=1.2em, itemsep=2pt]
  \item Students may say the period is 24 (the total span of data shown).
    Prompt: ``What is the \emph{smallest} input length before the pattern repeats?''
  \item Students may list the period as a range of months (``months 1--12'')
    rather than the length ($k = 12$). Distinguish: the period is a
    \emph{length}, not an interval.
\end{itemize}
\end{teachernote}

\vspace{0.08in}

%----------------------------------------------------------------------
% Tier A Key
%----------------------------------------------------------------------
\begin{tcolorbox}[colback=white, colframe=black!40,
    title=\textbf{Tier A --- Apply (Key)},
    fonttitle=\bfseries, arc=2mm, left=3mm, right=3mm, top=2mm, bottom=2mm]

\textbf{1.}\quad Table for two full periods:

\begin{center}
\renewcommand{\arraystretch}{1.5}
\begin{tabular}{|c|c|c|c|c|c|c|c|c|c|}
\hline
$t$ (sec) & 0 & 2 & 4 & 6 & 8 & 10 & 12 & 14 & 16 \\
\hline
$h(t)$ (ft) & \ans{0} & \ans{10} & \ans{0} & \ans{$-10$} & \ans{0} & \ans{10} & \ans{0} & \ans{$-10$} & \ans{0} \\
\hline
\end{tabular}
\end{center}

\vspace{6pt}
\textbf{2.}\quad \ans{Yes, the relationship is periodic. The sequence 0, 10, 0, $-$10, 0 repeats exactly every 8 seconds, over equal-length intervals.}

\vspace{4pt}
\textbf{3.}\quad $k =$ \ans{8} seconds.

\vspace{4pt}
\textbf{4.}\quad
\begin{itemize}[leftmargin=1.4em, itemsep=3pt]
  \item Increasing: \ans{$0 \leq t \leq 2$, $8 \leq t \leq 10$ (and repeats every 8 sec)}
  \item Decreasing: \ans{$2 \leq t \leq 6$, $10 \leq t \leq 14$ (and repeats every 8 sec)}
\end{itemize}

\vspace{4pt}
\textbf{5.}\quad Average rate of change from $t=0$ to $t=2$:

$\dfrac{h(2)-h(0)}{2-0} = \dfrac{10-0}{2} = $ \ans{5 ft/sec}

\end{tcolorbox}

\begin{teachernote}
\textbf{Common errors --- Tier A:}
\begin{itemize}[leftmargin=1.2em, itemsep=2pt]
  \item Students may list the interval $t=6$ to $t=8$ as increasing (it goes
    from $-10$ to $0$, which is indeed increasing). Accept this; it matches
    the formal definition.
  \item Some students divide by 8 instead of 2 when computing the average rate
    of change. Remind: divide by the \emph{width} of the interval
    ($b - a = 2 - 0 = 2$).
\end{itemize}
\end{teachernote}

\vspace{0.08in}

%----------------------------------------------------------------------
% Tier E Key
%----------------------------------------------------------------------
\begin{tcolorbox}[colback=white, colframe=black!40,
    title=\textbf{Tier E --- Extend (Key)},
    fonttitle=\bfseries, arc=2mm, left=3mm, right=3mm, top=2mm, bottom=2mm]

\begin{center}
\renewcommand{\arraystretch}{1.3}
\small
\begin{tabular}{|c|c|c|c|c|c|c|c|c|c|c|c|c|c|}
\hline
$t$ & 0 & 1 & 2 & 3 & 4 & 5 & 6 & 7 & 8 & 9 & 10 & 11 & 12 \\
\hline
$A(t)$ & 0 & 1.5 & 2 & 1.5 & 0 & $-1.5$ & $-2$ & $-1.5$ & 0 & 1.5 & 2 & 1.5 & 0 \\
\hline
$B(t)$ & 0 & 3 & 4 & 3 & 0 & $-3$ & $-4$ & $-3$ & 0 & 3 & 4 & 3 & 0 \\
\hline
\end{tabular}
\end{center}

\vspace{6pt}
\textbf{1.}\quad \ans{Both $A$ and $B$ are periodic. For $A$: the sequence 0, 1.5, 2, 1.5, 0, $-$1.5, $-$2, $-$1.5, 0 repeats with period $k = 8$ hours. For $B$: the same pattern of positions (scaled) repeats with period $k = 8$ hours.}

\vspace{4pt}
\textbf{2.}\quad Average rate of change of $A$, $t = 0$ to $t = 2$:

$\dfrac{A(2)-A(0)}{2-0} = \dfrac{2-0}{2} = $ \ans{1 m/hr}

\vspace{4pt}
\textbf{3.}\quad Average rate of change of $B$, $t = 0$ to $t = 2$:

$\dfrac{B(2)-B(0)}{2-0} = \dfrac{4-0}{2} = $ \ans{2 m/hr}

\vspace{4pt}
\textbf{4.}\quad \ans{Buoy B's wave height is changing more rapidly during $0 \leq t \leq 2$ (rate 2 m/hr vs.\ 1 m/hr for Buoy A). Physically, this means Buoy B's water surface is rising twice as fast as Buoy A's during this interval.}

\vspace{4pt}
\textbf{5.}\quad \ans{By EK 3.1.B.3, all characteristics of a periodic function repeat in every period. Since both $A$ and $B$ have period 8, the average rate of change on any interval $[t, t+2]$ equals the rate on $[t+8, t+10]$, $[t+16, t+18]$, and so on --- the same rate reappears at the same relative position in every cycle.}

\end{tcolorbox}

\begin{teachernote}
\textbf{Common errors --- Tier E:}
\begin{itemize}[leftmargin=1.2em, itemsep=2pt]
  \item Students may claim different periods for $A$ and $B$ because $B$ has
    larger outputs. The \emph{amplitude} differs but the period is the same
    (8 hours) for both --- reinforce that period is determined by the
    \emph{input} spacing, not the output magnitudes.
  \item Students answering \#5 should use the formal definition
    $f(x+k)=f(x)$ to justify the claim rather than just stating it.
\end{itemize}
\end{teachernote}

\end{document}
